% Supplementary Information for the QSS manuscript. Build: paper/manuscript_qss/build.sh (with main.tex).
% Numbers come from paper/FINDINGS.md; figures from paper/figures (make_figures.py, make_tau_sweep.py).
\documentclass[11pt]{article}
\usepackage[margin=1in]{geometry}
\usepackage[T1]{fontenc}
\usepackage[utf8]{inputenc}
\usepackage{mathptmx}
\usepackage{microtype}
\usepackage{amsmath,amssymb}
\usepackage{graphicx}
\usepackage{booktabs}
\usepackage{threeparttable}
\usepackage[font=small,labelfont=bf]{caption}
\usepackage{setspace}
\usepackage[round,authoryear]{natbib}
\usepackage{xcolor}
\usepackage{xr-hyper}
\usepackage[hidelinks]{hyperref}
\externaldocument{main}
\usepackage{lineno}
\linenumbers
\onehalfspacing
\graphicspath{{../figures/}}

\newcommand{\PiP}{\mathrm{PiP}}
\newcommand{\Peers}{\mathcal{P}}

\title{Supplementary Information for ``Percentile in Percentile: rating scholars against their career
cohort by the impact of their papers''}
\author{Panos Ipeirotis\thanks{Department of Technology, Operations, and Statistics, Stern School of Business, New York University.}\textsuperscript{,}\thanks{Corresponding author: \texttt{panos@stern.nyu.edu}}
  \and Sen Tian\thanks{Ciphero.}}
\date{}

\renewcommand{\thesection}{S\arabic{section}}
\renewcommand{\theequation}{S\arabic{equation}}

\begin{document}
\maketitle

\noindent This supplement gives the detailed checks that the main text summarizes. Section,
table, figure and equation numbers without an S refer to the main text.

\section{Factorization}
\label{supp:factor}
\noindent Main text: Section~\ref{sec:properties}.

PiP can be written as the author's productivity percentile times an impact multiplier. Let
$F(m) = |\{b \in \Peers : n_b(0) \le m\}| / |\Peers|$ be the productivity percentile. Unlike $G_\theta$, it counts
tied peers, so that it is positive for every member of the peer group. Defining the \emph{impact multiplier}
$Q_\tau(a) = \PiP_\tau(a) / F(n_a(0))$ gives
\begin{equation}
\PiP_\tau(a) = F\!\left(n_a(0)\right) \cdot Q_\tau(a).
\label{eq:factor}
\end{equation}
The identity holds by the definition of $Q_\tau$; it describes a score rather than explaining it. The multiplier
compares the author's standing on counts of articles above the cutoff with their standing on the number of articles,
and it exceeds one when the author ranks higher among peers on the former. Because $F$ counts ties and $G_\theta$ does
not, an author whose standing is the same at every threshold has $Q_\tau$ slightly below one.

In our data $Q_{1/2}$ exceeds one for 30--36\% of eligible authors, depending on career age. Across eligible authors
with a positive score who started in 1975--2005, at career ages 5, 10 and 20, the variance of $\log \PiP_{1/2}$ splits
into 0.74--0.82 from $\log Q$, about 0.30 from $\log F$ and a small negative covariance ($-0.04$ to $-0.11$). Most of
the variation in PiP across authors thus comes from the impact multiplier rather than from the number of articles. These
shares describe how the scores vary; they do not measure how much of PiP reflects research quality.

\section{Ties}
\label{supp:ties}
\noindent Main text: Section~\ref{sec:properties}.

The standing counts peers with strictly fewer articles. Mid-ranks instead count half of the tied
peers. They are a common choice for the percentile of a single paper \citep{bornmann2013use,waltman2013percentile},
but in the standing they would give credit to authors with nothing above a threshold: at high thresholds most peers
have no article, and an author tied with them at zero would stand at half their share. In the simulated cohort, an
author whose five articles all fall below the median would score 0.22 with mid-ranks instead of 0. The strict rule
keeps weak articles from earning credit through ties; its price is that the mean score of a peer group depends on how
common
ties are (Eq.~\ref{eq:mean}).

\section{The earlier, asymmetric form of the score}
\label{supp:asym}
\noindent Main text: Section~\ref{sec:relation}.

The earlier form of the score \citep{ipeirotis2024pip} is asymmetric; the website \url{https://www.pip-score.org} uses
it. The asymmetric form, as the website computes it, sorts an author's articles by percentile, $y_{(1)} \ge y_{(2)} \ge
\cdots$, and computes the area under the step curve that plots $y_{(k)}$ against the productivity percentile $F(k)$:
\begin{equation}
\text{PiP-AUC}(a) = \sum_{k=1}^{n_a(0)} y_{(k)}\,\bigl[F(k) - F(k-1)\bigr] = \int_0^1 F\!\left(n_a(\theta)\right)
d\theta .
\end{equation}
The second expression shows the difference from PiP: the asymmetric form compares the author's count at or above
$\theta$ with peers' \emph{total} counts rather than with their counts at or above $\theta$.

The asymmetric form is the product of $F(n_a(0))$ and a weighted mean of the author's own percentiles, which keeps it
bounded and monotone. But the weight of the $k$-th best article, $F(k) - F(k-1)$, is the share of peers with exactly $k$
articles, which is zero for every $k$ below the smallest article count among peers. Among the eligible authors of a
cohort, nobody has fewer than three articles, so an author's best two articles receive no weight at all; among Google
Scholar faculty profiles, the median author's best four articles receive none. This form is therefore dominated by
productivity: the productivity percentile accounts for 0.59--0.77 of the variance of its logarithm. Scored against the
simulated cohort, whose members all have at least three articles, an author with one article at the 99th percentile
scores zero, while an author with 40 weak articles scores 0.10. Applying a cutoff $\tau$ to this form causes a second
problem: it then compares partial counts with growing total counts, and its median at career age 10 halves between the
1960 and 2010 cohorts. The symmetric comparison in PiP removes both problems.

\section{The exact integral}
\label{supp:exact}
\noindent Main text: Section~\ref{sec:computation}.

To check that the classes do not drive the results, we also computed the exact integral for the 5.3 million eligible
authors at career age 10 (start years 1960--2015). Write $y_{(k)}$ for the author's $k$-th best percentile. On
$(y_{(k+1)}, y_{(k)}]$ the author has $k$ articles at or above $\theta$, and a peer has fewer exactly when its own
$k$-th best percentile is below $\theta$, so the integral reduces to sorting the peers' $k$-th best percentiles.
Against the exact score, the 40-class score is 0.010 higher on average at $\tau = 0.5$ (0.055 at most) and ranks
authors within cohorts almost identically: the rank correlation is 0.9994, and 96\% of each cohort's top 10\% is the
same (start years 1975--2010). At $\tau = 0.75$ and $0.9$ the correlation is 0.998 and 0.996 and the overlap 93\% and
85\%. None of the prediction results at career age 10 (Section~\ref{sec:prediction}) changes by more than 0.004 at any
of the three cutoffs, so we keep the 40 classes.

\section{Retrospective check of the cutoff}
\label{supp:holdout}
\noindent Main text: Section~\ref{sec:tau}.

The default was chosen from results that pooled all cohorts, so we
checked afterwards whether the choice depends on which cohorts it is made from. Because the check was made after the
default was chosen, it does not protect against specification search as an untouched holdout would. We split the
prediction settings that compare authors of one career age into a choice sample (start years 1975--1990) and a test
sample (start years 1991 and later). Stated as a rule, the default is the cutoff with the highest AUC for a future
top-1\% paper among cutoffs whose distribution is stable across cohorts ($\tau \le 0.6$). Applied to the choice sample
alone, the rule selects $\tau = 0.6$ at career age 10, with $\tau = 0.5$ within 0.001. The choice cohorts differ from
the test cohorts in two respects. The knee of their quantity--impact frontier (Figure~\ref{fig:frontier}) lies at a
slightly higher cutoff (between 0.6 and 0.75 rather than at 0.5). In those cohorts at career age 5, the top-10\% count
ties $\tau = 0.5$ on the top-1\% outcome and beats it on future citation impact, whereas $\tau = 0.6$ keeps its lead.

In the test cohorts both cutoffs keep every comparison reported in Sections~\ref{sec:prediction}--\ref{sec:early} at
career ages 5, 10 and 15: they beat the $h$-index on all three outcomes, no other indicator beats either of them on more
than one outcome, and both have the highest AUC for a future top-1\% paper. Across $\tau$ from 0.25 to 0.6 that AUC
varies by less than 0.01 in the test cohorts, so the exact choice in this range matters little for the outcome used to
make it.

\section{Replication across cohorts}
\label{supp:cohorts}
\noindent Main text: Section~\ref{sec:prediction}.

With more than a million authors per setting, sampling error is small: the
standard error of a Spearman correlation of 0.4 would be below 0.001 for independent observations. Authors share
papers, so the observations are not independent and this standard error is optimistic. The more relevant question is
whether the differences hold in different populations. We repeated the age-10 comparison separately for each start
year from 1975 to 2000, 26 cohorts of 17,867 to 122,114 authors. PiP beats the $h$-index on all three outcomes in all
26 cohorts, by 0.06--0.10 on the number of future top papers. It beats the top-10\% count on the number of future top
papers (by 0.05--0.08) and on the top-1\% outcome in all 26 cohorts, and falls behind it on future citation impact in
20 (by up to 0.07). Against I3 and the article count, PiP wins on the top-1\% outcome and on future citation impact in
every cohort, and on the number of future top papers in 12 and 11 of the 26, with differences between $-0.03$ and
$+0.04$.

\section{Early detection under the three-article rule}
\label{supp:min3early}
\noindent Main text: Section~\ref{sec:early}.

Requiring an $h$-index of three by career age 3 keeps only 21\% of the
authors who have three articles by then, and it selects the population by one of the indicators being compared. We
therefore repeated the analysis for all 1,455,891 authors with at least three articles at career age 3, with peers
defined by the same rule. PiP still leads the $h$-index at every measurement age: the top 10\% by PiP contains 24\%
future stars at age 3 against 20\% for the $h$-index, and 31\% against 27\% at age 5. PiP reaches the $h$-index's age-10
accuracy 1.2 to 1.7 years earlier. It also stays ahead of the indicators in Table~\ref{tab:early} on precision and on
the AUC for a future top-1\% paper at every age, but I3 matches or exceeds it on the AUC for future stars from about
career age 6.

\section{Agreement on the Google Scholar profiles}
\label{supp:gsagree}
\noindent Main text: Section~\ref{sec:honors}.

A last check measures agreement rather than prediction. On the Google Scholar profiles, where scores and honors are
measured at the same time, the indicators are close: PiP is slightly ahead of the $h$-index for holding an honor
($+0.004$ [0.001, 0.008]) and for having a Wikipedia article ($+0.012$ [0.009, 0.015]). Total citations and the
top-10\% count do as well for holding an honor (their paired differences from PiP include zero) and slightly better
for having a Wikipedia article (by 0.005 and 0.004).

\section{Author disambiguation}
\label{supp:disamb}
\noindent Main text: Section~\ref{sec:robustness}.

Semantic Scholar's author records come from automatic disambiguation, which can split
one person across several records or merge different people. To check that the results are not artifacts of this
algorithm, we use the author-curated Google Scholar profiles. We matched the 16,218 profiles to Semantic Scholar through
their publications. A profile's work matches a Semantic Scholar paper on its normalized title (at least 25 characters)
and a publication year within one year. The profile's counterpart is the Semantic Scholar author record with the same
surname that holds the most matched works. We find a confident counterpart (at least five matched works and at least
30\% of the profile's titles) for 14,511 profiles (89\%). For the median person, 76\% of their Google Scholar works are
found in Semantic Scholar, and 90\% of those found sit under the counterpart. Splitting is common, though: 64\% of
people have at least two of their works under a second record with the same surname and first initial, and 42\% under a
third. The counterpart is smaller than the profile (a median 67\% of the profile's cited works, and 78\% of its
$h$-index), and its start year is later than the profile's for 47\% of people, because early work is missing or sits
under another record.

Despite these differences, the two sources rank people similarly. For the 9,796 people who have scores in both sources
(Semantic Scholar start year 1995 or later, evaluated in 2025), the Spearman correlation between sources is 0.82 for the
$h$-index and 0.78 for PiP. Each source computes percentiles against its own population: all eligible Semantic Scholar
authors, and the Google Scholar profiles. For PiP, the agreement is lower for people split across three or more records
(0.73) than for people with at most one record (0.76), but not by much. We read these comparisons as evidence that the
results above are not artifacts of one disambiguation algorithm, and as a warning that individual Semantic Scholar
scores can be off when a person's work is split.

\section{Survivorship}
\label{supp:surv}
\noindent Main text: Section~\ref{sec:robustness}.

The population includes every eligible author, not only those who stayed in research, and
the prediction results of Section~\ref{sec:prediction} count authors who stop publishing as having no future top
papers. This choice favors indicators that predict survival, that is, whether an author continues to publish. Measured
at career age 5, the indicators predict whether an author is still publishing at career ages 11--15 with an AUC of
0.68 for the article count, 0.62 for PiP, 0.58 for the $h$-index and 0.54 for citations. Continuing to publish is a
matter of volume, so the indicators that weigh volume most predict it best. PiP still leads the article count, the
best predictor of survival, on the top-1\% outcome and on future citation impact.

\bibliographystyle{apalike}
\bibliography{references}

\end{document}
